% Generado por sistema/backend ReporteTex.java -- no editar a mano.
% Datos: Base de datos: 148 importado (titulados) / Base de datos: 53 importado (empleadores)
\subsubsection*{C2. Sector $\times$ pertinencia del puesto}

\textit{Propósito.} Determinar si la concordancia entre la titulación y el empleo depende del sector, público o privado, en que trabaja el titulado. Respalda el criterio ARCU-SUR 3.2.3. Se aplica sobre: trabajan en una organización, con nivel de significancia $\alpha = 0.05$.

\textit{Paso 1. Contar con la muestra: tabla de contingencia} (frecuencia observada $FO_{ij}$, $N = 108$).

\begin{mitabla}{C2: frecuencia observada $FO_{ij}$}
\begin{tabular}{lcc|c}
\toprule
Sector & En desacuerdo & De acuerdo & Total fila \\
\midrule
Público & 2 & 19 & 21 \\
Privado & 8 & 79 & 87 \\
\midrule
Total columna & 10 & 98 & $N=108$ \\
\bottomrule
\end{tabular}
\end{mitabla}

\textit{Paso 2. Calcular la frecuencia esperada} de cada celda, $FE_{ij} = (\text{total fila } i)(\text{total columna } j)/N$. Por ejemplo, para la celda ``Público -- En desacuerdo'': $FE_{11} = 21 \times 10 / 108 = 1.9444$.

\begin{mitabla}{C2: frecuencia esperada $FE_{ij}$}
\begin{tabular}{lcc}
\toprule
Sector & En desacuerdo & De acuerdo \\
\midrule
Público & $21\cdot10/108 = 1.94$ & $21\cdot98/108 = 19.06$ \\
Privado & $87\cdot10/108 = 8.06$ & $87\cdot98/108 = 78.94$ \\
\bottomrule
\end{tabular}
\end{mitabla}

\textit{Paso 3. Calcular el estadístico chi-cuadrado,} sumando $(FO_{ij}-FE_{ij})^2/FE_{ij}$ de cada celda.

\begin{mitabla}{C2: cálculo de $X_0^2$ celda por celda}
\begin{tabular}{L{5.4cm}cccc}
\toprule
Celda & $FO_{ij}$ & $FE_{ij}$ & $(FO_{ij}-FE_{ij})^2$ & $(FO_{ij}-FE_{ij})^2/FE_{ij}$ \\
\midrule
Público -- En desacuerdo & 2 & 1.9444 & 0.0031 & 0.0016 \\
Público -- De acuerdo & 19 & 19.0556 & 0.0031 & 0.0002 \\
Privado -- En desacuerdo & 8 & 8.0556 & 0.0031 & 0.0004 \\
Privado -- De acuerdo & 79 & 78.9444 & 0.0031 & 0.0000 \\
\midrule
\multicolumn{4}{r}{$X_0^2 = \sum$} & 0.0022 \\
\bottomrule
\end{tabular}
\end{mitabla}

\textit{Paso 4. Determinar el estadístico teórico.} Con $k = 2$ filas y $m = 2$ columnas, los grados de libertad son $(k-1)(m-1) = (2-1)(2-1) = 1$. De la tabla de chi-cuadrada: $X_{0.05,\,1}^2 = 3.841$.

\textit{Paso 5. Establecer la relación entre estadísticos:} $X_0^2 < X_{0.05,\,1}^2 \;\Rightarrow\; 0.0022 < 3.841$.

\textit{Paso 6. Conclusión.}

\begin{center}
\begin{tabular}{l}
Si ($X_0^2 < X_{0.05,\,1}^2$)\\
\quad \fbox{Independencia}\\
Sino\\
\quad Relación entre las variables\\
FinSi
\end{tabular}
\end{center}

No se puede rechazar que sector y pertinencia del puesto sean independientes.

\textit{Observación:} 1 de 4 celdas tienen $FE < 5$, por lo que la prueba pierde confiabilidad y la conclusión debe tomarse con cautela.

